\documentclass[10pt]{article}
\usepackage[margin=20mm]{geometry}
\usepackage{amsmath,amssymb,fontspec,unicode-math,xurl,hyperref}
\IfFontExistsTF{DejaVu Serif}{\setmainfont{DejaVu Serif}}{\setmainfont{Latin Modern Roman}}
\setmathfont{latinmodern-math.otf}
\setlength{\parindent}{0pt}\setlength{\parskip}{6pt}\sloppy
\begin{document}
ELVIA Scientific Verification Report 0075
MP-Nav: Enhancing Data Poisoning Attacks against Multimodal Learning
Paper I25-43444 | Supported outcome: Minor: bounded pairing inconsistency
Review date: 24 October 2026 | Human scientific decision pending\par
\section*{1. Identity, version and inspected material}
Jingfeng Zhang, Prashanth Krishnamurthy, Naman Patel, Anthony Tzes and Farshad Khorrami. ICML 2025, Proceedings of Machine Learning Research 267:75283-75296 [S1]. The reviewed object is the 14-page publisher-linked PDF preserved as `papers/0075/0075.pdf`. All primary-paper locations use one-based PDF pages, identical to printed internal page numbers. Its SHA-256 was verified against the source-acquisition record, and the original bytes were preserved.\par
The main text (pp1-8), Appendix A and extended experiments (pp12-14), all six tables and the main experimental figures were inspected. Selected baseline evidence was checked in Yang et al.’s ICML 2023 AtoB paper [S2] and Xu et al.’s NeurIPS 2024 ShadowCast paper [S3]. Five small files from the official ShadowCast repository were inspected statically at commit `b13f877a2580de6a6a50833017b699dee87353ce`, dated 15 January 2025 [S4]. They corroborate candidate-pool construction versus injection. They are baseline code, not a recovered MP-Nav implementation.\par
No attributable MP-Nav implementation was located through the publisher metadata, first-author research page/public repositories and bounded exact-title searches. The publisher lists PDF and OpenReview links and an empty software-extras field. OpenReview access returned a browser-verification page or HTTP 403, so discussion and revision history were not read. No separate arXiv version or scientific erratum was located in the inspected records. The publisher PDF path history returned one October 2025 commit; its generic metadata-fix message is not a scientific correction. These are bounded retrieval results, not proof that additional material does not exist.\par
\section*{2. Outcome, definitions and limits}
One confirmed concern supports Minor. P1 is an ordered image/text pairing error in Algorithm 1’s AtoB branch, inconsistent with the stipulated attack direction. The predecessor report’s P2 allegation is not established: selecting a subset from a candidate poison pool does not show that the learner’s actual training set contained unpoisoned examples of the same concept. The baseline-composition question is retained as Needs Evidence. No finding establishes that MP-Nav’s actual training used the wrong pairing or that its reported gains are false.\par
Throughout the paper x denotes an image and y a text caption. O and T are distinct original and target concepts. In text-to-image retrieval, O is the query-text concept and T the image concept to be incorrectly retrieved. In the VQA setting, the original input is an image from O and the desired output mentions T. The two tasks therefore have different source modalities. ShadowCast constructs an image visually resembling T but with features resembling O, and pairs it with text about T.\par
MP-Nav selects concept pairs using cosine similarity of mean image/text embeddings, then ranks instances by cosine proximity to a concept center built from a weighted image/text representation. The paper presents this as an empirical selection strategy. Its p7 limitation explicitly acknowledges that neither component always improves an attack. No universal robustness or optimality theorem was imposed in this review.\par
All numerical checks here recompute displayed summaries, identities or finite comparisons. No model was trained, no poison was generated, no author code was executed, no model/API experiment was paid for, and no raw seed-level result was replayed. Human scientific disposition remains pending.\par
\section*{3. Problem 1 — Algorithm 1 reverses the stipulated AtoB pairing}
Source claim and exact conflict. Section 2.2, p3, Section 2.2.1, p4, and the Hit@K/MinRank definitions on p6 consistently specify a text query from O retrieving an image of T. Section 2.2.1 consequently pairs target images x\_T with original captions y\_O. Algorithm 1 on the same p4 instead pairs x\_O with y\_T and replaces an original same-concept pair with that pair [S1]. The cited AtoB paper’s Section 3.2, p3, also defines its single-target-label attack by combining original-class texts with target-class images [S2].\par
Complete finite identity check. Choose distinct concepts O=bear and T=lion, with correctly labelled bear and lion images and corresponding captions. No ambiguous or multi-label image is needed. The body recipe produces (lion image, bear caption); the algorithm produces (bear image, lion caption). These are different ordered pairs and associate opposite text-query concepts with image concepts. The discrepancy exists before any optimizer or training randomness enters. The supported statement concerns sample construction; it is not a claim that fitting either pair must guarantee a particular learned retrieval output.\par
Alternative readings and counterevidence. A local notation error is plausible because the body gives the intended pairing clearly. A global exchange of O and T would break the objective and metric definitions, and would not preserve Algorithm 1’s separately coherent ShadowCast branch. Cosine similarity being symmetric in its two vector arguments does not make the two cross-modal concept assignments identical: swapping concepts also changes which encoder supplies each vector. For example, unit image embeddings (1,0),(0,1) and text embeddings (4/5,3/5),(0,1) give S[O,T]=0 and S[T,O]=3/5 [C1]. This illustrates the distinction; it does not prove failure of the concept-selection heuristic.\par
Algorithm 1’s word ’replace’ may refer to the attacker’s own candidate dataset before injection. It therefore does not independently establish unauthorized removal of clean victim data. No implementation was obtained that would resolve which pairing was actually used in the reported runs, and the table labels alone cannot establish that the experiments ran in reverse.\par
Impact and precise correction. The executable recipe a reader would reconstruct differs from the stated TIR objective; this is a substantive method-specification error, not merely spelling. The supported scope is nevertheless local specification, so the judgment is Minor. Replacing the AtoB branch’s output pair by (x\_T,y\_O), with its corresponding source/target selection indices stated consistently, restores the ordered-pair agreement with Section 2.2.1. This logical consistency repair is verified; training behavior, budget indexing, the chosen similarity orientation and author adoption remain unverified.\par
\section*{4. Baseline clarification - candidate pool and learner-training composition}
Accurate claim and location. The introduction on p2 says that the prior example required poisoning all instances of the Biden concept in the training data. Section4.2, p7 additionally states that perturbations were applied to all277 instances and that the learner was assumed to have no benign Biden examples. The phrase about perturbing all277 can reasonably describe candidate-pool generation, while the no-benign statement concerns the composition of the learner’s final training set. These are different propositions [S1].\par
What the prior evidence establishes. Xu et al., Section4.1 p6 distinguishes a pool of200–300 crafted candidates from the \(M\) samples randomly injected into the clean instruction-tuning set, with \(M\) in5,10,20,30,50,100,150,200. Section4.2 p7 reports strong Trump-to-Biden results at budgets around50 and above [S3]. The pinned preparation code samples num\_poison distinct candidate indices and appends only those entries to the preexisting clean-data annotations; it does not append the unused candidate images as clean same-concept examples [S4]. These source and code facts remain valid.\par
Why the predecessor counterargument fails. If the preexisting clean instruction set contains no Biden examples, all Biden examples in the learner’s final training set can be poisoned even when only50 out of277 candidates are injected. Thus \(M<277\) is compatible with the paper’s all-in-training/no-benign reading. Neither the baseline paper’s budget sweep nor its preparation code, as inspected, demonstrates successful attacks with retained clean Biden examples in the learner’s training set. The exact candidate count277 and the actual composition of the clean instruction data were not independently verified. An acceptance search of the retained baseline text found no direct clean-Biden mixture result resolving this issue; that bounded check does not prove such evidence cannot exist.\par
Disposition and scope. The stronger assertion that success requires absence of clean same-concept examples is not established by the inspected baseline budget sweep, but the saved evidence also does not refute it under the operative training-composition condition. The predecessor’s confirmed-Minor mischaracterization finding is therefore withdrawn and retained as Needs Evidence, with the candidate-versus-injected distinction recorded as a clarification. MP-Nav’s explicitly added-clean-Biden experiment remains a distinct comparison. Resolution would require the actual baseline clean-data concept composition or a matched experiment retaining benign destination-concept examples. No author-code execution, dataset download or new attack experiment was performed.\par
\section*{5. Numerical and methodological checks that constrain the interpretation}
Table 2, p6 is internally consistent with its reported improvement. The combined configuration changes boat-to-dog into boat-to-kit and increases displayed Hit@1 by 0.038, Hit@5 by 0.204 and Hit@10 by 0.172 (17.2 percentage points, 34.06\% relative to 0.505); displayed MinRank decreases by 1.103. Holding boat-to-dog fixed, instance selection improves Hit@10 by 0.113 and reduces MinRank by 0.753. Holding boat-to-kit fixed, adding instance selection improves Hit@10 by 0.067 and reduces MinRank by 0.413. These calculations are differences of reported medians, not paired-seed treatment estimates.\par
Table 5, p14 contains nine source-concept groups and four configurations per group. An independent parser extracted all 36 rows, with source-line references and displayed deviations, and counted changes in the favorable direction [C1]. The following counts are descriptive comparisons of central values, not significance tests:\par
Comparison | Concept pair | Improved Hit@1 /5 /10 | Improved MinRank\par
Instance vs baseline | Same pair | 6/9 ; 7/9 ; 7/9 | 6/9\par
Concept vs baseline | Changed target | 8/9 ; 8/9 ; 8/9 | 7/9\par
Combined vs baseline | Changed target | 8/9 ; 8/9 ; 8/9 | 7/9\par
Combined vs concept | Same selected pair | 5/9 ; 5/9 ; 6/9 | 6/9\par
\section*{5.1. Author-disclosed counterevidence and causal scope}
The table shows meaningful setbacks. For skateboard, increasing the selected pair’s similarity from 0.286 to 0.293 changes Hit@10 from 0.581 to 0.041 for concept-only selection and worsens MinRank from 12.126 to 68.315. For sheep-to-snowboard, instance selection changes Hit@10 from 0.606 to 0.555 and worsens all four central metrics. For banana-to-skis, Hit@10 improves from 0.425 to 0.446 while MinRank worsens from 25.369 to 41.360; this is possible because the metrics summarize different parts of the rank distribution. The paper explicitly discusses non-guaranteed improvements on p7, so these are preserved counterevidence rather than newly discovered violations of a universal theorem.\par
Choosing an easier target is allowed by the stated threat model. Comparing absolute post-poison success for different targets therefore addresses a legitimate practical objective. It does not separately estimate the increment caused by poisoning unless matched clean results for each concept pair are supplied. For example, a hypothetical baseline pair with clean/poison Hit@10 of 0.10/0.50 has increment0.40, while a selected pair at0.60/0.68 has higher final success but increment0.08. Both readings can coexist; these illustrative numbers are not observations from the paper. The study’s fixed-pair instance comparisons provide additional evidence that the changed-target contrast alone would not.\par
An at-least-one-target Hit@K metric can also depend on target-gallery prevalence if counts differ, but an imbalance was not established here. The text’s aggregate dataset counts do not justify asserting unequal target counts. Raw per-pair gallery membership and clean attack metrics would resolve that question. Thus causal attribution of a poison-induced increment is Needs Evidence; the observed post-poison improvement itself is not disproved.\par
Uncertainty and utility. Table 2 and Figure 2 report medians with standard deviations over five training seeds; Figure 4 also describes five training trials. Figure 3 instead reports mean and standard deviation across three evaluation prompts. These are not interchangeable with confidence intervals or repeated-training uncertainty. No unreported formal significance test was presumed. Tables 3/6 show retrieval central values within 0.2 percentage points below their corresponding clean values at worst; Table 4’s largest displayed GQA reduction is 59.5 to 59.0 for the combined method at the reported 0.27\% rate. These support comparable aggregate utility on the listed benchmarks, not a proved equivalence margin or preservation of every class.\par
Remaining protocol limits. Lambda in the mixed representation, detailed tie/rounding behavior in top-budget selection, exact Food101 poison-rate denominator/subset construction and the actually used TIR optimizer configuration were not independently recoverable from a MP-Nav implementation. The printed TIR initial learning rate 0.2 also differs from the cited AtoB paper’s 10\textasciicircum{}-5; this is a reproducibility question, not proof that either reported experiment is false. The p6 reference to COCO having 25 images per concept occurs in the Flickr-PASCAL endpoint explanation; the described PASCAL split provides the relevant 25-per-class context, so no new numerical failure is inferred from that label.\par
\section*{6. Dependence, repair status and remaining work}
The supported severity is Minor based on P1. The empirical proposal to examine semantic similarity and central instances remains testable and has favorable as well as unfavorable reported comparisons. P1 affects the written AtoB construction; the predecessor’s P2 counterargument confused the candidate pool with learner-training composition and is not a confirmed defect. The existing evidence establishes no failure of the core selection inference or published experimental direction. These scoped conclusions do not require the entire method to fail before a core error could count as Major.\par
Resolving actual attack direction requires the MP-Nav data-construction implementation or saved training annotations. Resolving the size and causal source of its gains requires exact splits, gallery labels, clean per-pair metrics, seed-level results and a faithful training reproduction. Neither absence of code nor an unresolved subitem is used to raise the paper’s supported severity.\par
Item / judgment | Repair verification | Adoption / remaining work\par
P1: Minor pairing inconsistency | Ordered-pair correction matches body objective analytically | Not adopted; actual training pairs and indexing unresolved\par
P2: Needs Evidence / pool clarification | Subset injection does not prove retained clean same-concept training examples | Need actual baseline clean-data composition or a matched mixture result\par
Selection limitations: clarification | Regressions and author qualification retained; Table5 counts recomputed | No universal guarantee or significance certificate asserted\par
Poison-increment attribution: Needs Evidence | Changed-target outcomes do not identify matched clean-to-poison increments | Per-pair clean and seed-level data needed\par
Overall: Minor | No author code execution or empirical repair validation | Human scientific decision pending\par
\section*{7. Consulted references and evidence identifiers}
[S1] Zhang et al., ICML 2025/PMLR267. https://proceedings.mlr.press/v267/zhang25am.html . Exact primary PDF: https://raw.githubusercontent.com/mlresearch/v267/main/assets/zhang25am/zhang25am.pdf . Local source: `papers/0075/0075.pdf`; extraction and selected page images: `sources/0075/`. Publisher metadata and PDF-path history are preserved in publisher\_metadata.md and publisher\_pdf\_history.json.\par
[S2] Yang et al., Data Poisoning Attacks Against Multimodal Encoders, ICML 2023/PMLR202:39299-39313, publisher version. Section 3.2 p3 and experimental setup p4 inspected. https://proceedings.mlr.press/v202/yang23f.html ; PDF https://proceedings.mlr.press/v202/yang23f/yang23f.pdf . Saved as `sources/0075/atob2023.pdf` and atob2023.txt. This is selected baseline verification, not a full review of that paper.\par
[S3] Xu et al., Shadowcast: Stealthy Data Poisoning Attacks against Vision-Language Models, NeurIPS 2024 conference PDF, 32 pages. Sections 4.1-4.2 pp6-7 and related budget material inspected. https://proceedings.neurips.cc/paper\_files/paper/2024/file/6a2e30664b9647f97d7b9275358d083c-Paper-Conference.pdf . Saved as `sources/0075/shadowcast2024.pdf` and shadowcast2024.txt. Its SHA-256 is `9fc0b72c4a13a71cd553ce3f5ce2bee2ce3159219198d67e8a98b1289a97b528`.\par
[S4] Official baseline code, umd-huang-lab/VLM-Poisoning, commit `b13f877a2580de6a6a50833017b699dee87353ce`. https://github.com/umd-huang-lab/VLM-Poisoning/tree/b13f877a2580de6a6a50833017b699dee87353ce . The five preserved files retain their paths under `sources/0075/shadowcast\_repository/`. Commit/tree metadata and per-file acquisition hashes are separately preserved. The README also explains that Biden\_base\_Trump\_target means the paper’s Trump-to-Biden attack, distinguishing construction terminology from input/output direction.\par
[S5] First-author research page, inspected for an MP-Nav software/version link: https://zjfheart.github.io/research/data-safety-governance.html . OpenReview identifier zy7VeNtSLM; direct forum/PDF retrieval was blocked as recorded in support\_acquisition.json.\par
[C1] Independently written `verification/0075/independent\_checks.py`; results independent\_checks.json and table5\_extracted.json. These contain exact decimal differences, nine-group comparison counts, the ordered-pair identity check and the unit-vector cross-modal asymmetry illustration.\par
\section*{Supporting checks}
The preserved primary PDF is 664777 bytes, SHA-256 `e01538b863f8becb6a8cd54a9f75d7713c785c396cad0e98af1c0bc01fc1f3c4`. source\_verification.json records receipt matching. Existing Poppler utilities extracted UTF-8 layout text and rendered selected pages; source pp4,7,8 and14 were visually checked for algorithm, qualification, utility and table evidence.\par
It uses only standard-library parsing, Decimal arithmetic and data identities; it imports no author module. It completed successfully and parsed 36 Table 5 rows. The ordered-pair example and similarity illustration verify a specification distinction, not a learned-model prediction. No statistical significance or causal effect is inferred from the descriptive counts.\par
Actual retrieval limitations were preserved: OpenReview forum retrieval returned a challenge page; PDF retrieval returned HTTP 403, and a subsequent file check confirmed that no OpenReview PDF had been saved. An initial guessed raw-GitHub route for the older AtoB PDF returned 404; the publisher’s actual PDF URL then succeeded. No failed or timed-out mutation was blindly repeated. No MP-Nav code, training checkpoints or raw experiment logs were obtained; baseline image data were not downloaded.\par
Source and verification files remain in their frozen input directories.\par
Human scientific decisions remain pending, and author adoption has not been requested.\par
\end{document}
